\section{Увод}
\label{sec:uvod}

Квалифицираният електронен подпис (КЕП) отдавна не е техническа рядкост. С него се
подават заявления до администрацията и документи до съда, сключват се договори и се
подписват трудови документи. Удостоверение за КЕП вече може да се получи и без лична среща
с доставчика: правото на ЕС допуска самоличността на лицето да бъде проверена и без
физическото му присъствие, с методи, които гарантират идентифицирането му „с висока степен
на надеждност“ (чл.~24, пар.~1а, б.~„в“ eIDAS).\footnote{Регламент (ЕС) №~910/2014 на
Европейския парламент и на Съвета от 23 юли 2014~г. относно електронната идентификация и
удостоверителните услуги при електронни трансакции на вътрешния пазар и за отмяна на
Директива 1999/93/ЕО, ОВ L~257, 28.08.2014~г., с.~73 (eIDAS). Регламентът е изменен
съществено с Регламент (ЕС) 2024/1183 и се цитира по консолидирания текст към
18.10.2024~г.} Това улеснява гражданския оборот, но поражда и нов вид спор. Едната страна
представя документ с КЕП, а лицето, посочено като подписал, твърди, че никога не го е
подписвало.

Отговорът на въпроса кой трябва да докаже авторството на такъв документ се намира на три
равнища. Регламентът приравнява КЕП на саморъчния подпис (чл.~25, пар.~2 eIDAS). Законът
за електронния документ и електронните удостоверителни услуги (ЗЕДЕУУ) разграничава автора
и титуляря на изявлението и ограничава оспорването на авторството (чл.~4 и~18а ЗЕДЕУУ).
Гражданският процесуален кодекс (ГПК) определя как се представя електронният документ,
кой носи доказателствената тежест при оспорването му и как се проверява истинността му
(чл.~184 и~193--194 ГПК). Трите акта са приети по различно време и с различна логика,
затова съчетаването им не е очевидно.

Темата е актуална и защото правото на ЕС в тази област се променя. Регламент (ЕС)
2024/1183 въведе изчерпателен списък на методите за проверка на самоличността при
издаването на квалифицирани удостоверения и уреди като отделна квалифицирана услуга
управлението на устройства за създаване на КЕП от разстояние (чл.~24, пар.~1а и чл.~29а
eIDAS). Референтните стандарти за проверката на самоличността бяха приети през 2025~г. и
ще се прилагат от 19~август 2027~г. Същевременно ЗЕДЕУУ не е изменян от 2019~г.

\textit{Обект} на изследването са отношенията между титуляря на КЕП, адресата на
подписания документ и доставчика на удостоверителни услуги, а \textit{предмет}~--
доказателствената сила на документа с КЕП в гражданския процес и правилата за
оспорването му. \textit{Целта} е да се изясни кой и как трябва да докаже, че документ с
КЕП не изхожда от лицето, посочено като негов автор. \textit{Задачите} са: да се определи
правното действие на КЕП; да се покаже какво доказва подписът от техническа гледна точка;
да се анализират представянето и доказателствената сила на документа по ГПК; да се
изяснят редът, доказателствената тежест и средствата за проверка при оспорване; да се
разгледат злоупотребите и разпределението на риска между титуляря, адресата и доставчика;
да се формулират изводи и предложения.

Използвани са нормативният метод, систематичното тълкуване и сравнението между отделните
видове електронни подписи. Анализирани са практиката на Върховния касационен съд (ВКС) и
на Съда на Европейския съюз (СЕС), както и казус от административнонаказателната практика
на Комисията за регулиране на съобщенията (КРС), разгледан в курса. Работата е ограничена
до гражданския процес. Наказателноправните последици от неистинността на документа са
засегнати само доколкото ГПК препраща към тях.
